\newpage
\section{问题分析}

\subsection{问题一分析}
交通网络的图结构构建需要\textbf{邻接矩阵}定义。距离阈值法（固定半径内相连）和K近邻法（最近K个节点相连）各有优劣：前者保持局部拓扑，后者保证节点连通性。我们结合两者优点，构建混合邻接策略。

\subsection{问题二分析}
短时交通流预测的关键是\textbf{时空特征提取}。GraphSAGE通过聚合邻居节点信息学习空间特征，LSTM通过记忆单元捕捉时间依赖。两层图卷积聚合直接和间接邻居信息，充分捕捉交通流的扩散效应。

\subsection{问题三分析}
多路口协同控制的难点在于\textbf{状态空间爆炸}。N个路口的组合状态数为$|S|^N$，当N较大时难以枚举。我们采用分布式控制架构：每个路口作为独立Agent，通过消息传递协调相邻路口。
